\subsection{Welfare Decomposition and Counterfactuals}

This section decomposes the welfare effects of stablecoins into their underlying channels. We conduct a set of counterfactual exercises that isolate the cost and information mechanisms.

\subsubsection{Counterfactual 1: No Information Channel}

In this exercise, we shut down the information channel by fixing the precision of the public signal:
\begin{equation}
\sigma_s^2(q_s) = \bar{\sigma}_s^2.
\end{equation}

Stablecoins retain their cost advantage, but no longer improve information.

\textbf{Result.} Welfare increases monotonically with stablecoin usage.

\textit{Interpretation.} In the absence of the information channel, stablecoins are purely efficiency-enhancing. There is no coordination effect, and therefore no increase in fragility.

\subsubsection{Counterfactual 2: No Cost Advantage}

In this exercise, we eliminate the transaction cost advantage of stablecoins:
\begin{equation}
C_s(q_s) = C_c(q_s).
\end{equation}

Stablecoins retain their informational role but are no cheaper than cash.

\textbf{Result.} Welfare decreases with stablecoin usage.

\textit{Interpretation.} Without cost advantages, stablecoins provide no efficiency gains. Their only effect is to improve coordination, increasing crisis probability and reducing welfare.

\subsubsection{Counterfactual 3: Full Model}

In the baseline model, both channels operate simultaneously.

\textbf{Result.} Welfare is non-monotonic in stablecoin usage and exchange rate misalignment.

\textit{Interpretation.} The full model highlights the interaction between efficiency and fragility. At low levels of misalignment, efficiency gains dominate. At high levels, fragility costs dominate.

\subsubsection{Decomposition}

We summarize the welfare effect as:
\begin{equation}
\Delta W = \underbrace{\Delta W^{cost}}_{>0} - \underbrace{\Delta W^{info}}_{>0}.
\end{equation}

The cost channel always improves welfare, while the information channel always increases fragility. The net effect depends on the level of misalignment.

\subsubsection{Policy Implications}

The decomposition yields several policy-relevant insights:

\begin{itemize}
\item In low-misalignment environments, stablecoins are welfare-improving
\item In high-misalignment environments, stablecoins can be destabilizing
\item Policies that limit coordination (e.g., transaction frictions or regulation) may mitigate fragility
\end{itemize}

Thus, the impact of stablecoins depends critically on the underlying macroeconomic environment.
